\documentclass{covermaker}
\doctitle{Architecture \& Design Document}
\docsubtitle{Structure, decisions, compliance and model evaluation (arc42)}
\docid{CM-ADD}
\docstatus{Draft}
\input{glossary}

\begin{document}
\cmtitlepage

\begin{revisionhistory}
  0.1 & 2026-11-03 & Initial architecture, compliance analysis and model research (M0). \\
\end{revisionhistory}

\tableofcontents

\chapter{Introduction and goals}
\CM{} turns local audio files into selectable \glspl{stem} and
\glspl{backingtrack}. The main quality goals are:
\begin{enumerate}
  \item \textbf{Compliance by design}: no acquisition of audio from streaming
        platforms; transparent model licences.
  \item \textbf{Robustness}: long batch jobs survive crashes, out-of-memory
        errors and restarts.
  \item \textbf{Adaptivity}: good results from CPU-only machines up to large
        GPUs, without disturbing a running \gls{daw}.
  \item \textbf{Lightness}: a small installer and a responsive interface.
\end{enumerate}

\chapter{Context and scope}
\begin{center}
\begin{tikzpicture}[
  node distance=1.4cm and 2.2cm,
  box/.style={draw=cmaccent,fill=cmsoft,rounded corners,align=center,minimum width=3.2cm,minimum height=1cm},
  ext/.style={draw=cmmuted,rounded corners,align=center,minimum width=3.2cm,minimum height=1cm},
  arr/.style={-{Stealth[length=2.5mm]},thick,cmmuted}]
  \node[box] (app) {\textbf{\CM{}}\\desktop app};
  \node[ext,left=of app] (user) {User\\(musician, producer)};
  \node[ext,above right=0.6cm and 2.2cm of app] (models) {Model hosts\\(original sources)};
  \node[ext,right=of app] (cc) {Creative Commons\\catalogues};
  \node[ext,below right=0.6cm and 2.2cm of app] (gh) {GitHub Releases\\(updates)};
  \node[ext,below=of app] (fs) {Local files\\and folders};
  \draw[arr] (user) -- node[above,font=\small]{uses} (app);
  \draw[arr] (app) -- node[sloped,above,font=\small]{download weights} (models);
  \draw[arr] (app) -- node[above,font=\small]{licensed items} (cc);
  \draw[arr] (app) -- node[sloped,above,font=\small]{updates} (gh);
  \draw[arr,{Stealth[length=2.5mm]}-{Stealth[length=2.5mm]}] (fs) -- node[left,font=\small]{audio in / stems out} (app);
\end{tikzpicture}
\end{center}

All audio processing happens on the user's machine. The only outbound
connections are to model hosts chosen by the user, to optional Creative Commons
catalogues and to GitHub Releases for updates. No streaming platform is
contacted.

\chapter{Solution strategy}
\begin{tabularx}{\linewidth}{@{}lY@{}}
\toprule
\textbf{Concern} & \textbf{Approach} \\
\midrule
Desktop shell & Tauri 2 (Rust) with WebView2; Rust only handles operating-system integration. \\
User interface & SvelteKit (static adapter) + TypeScript, waveforms with wavesurfer.js and Web Audio. \\
Engine & Python 3.12 package managed by uv; FastAPI + WebSocket on 127.0.0.1 with a per-session token. \\
Separation & \texttt{audio-separator} models in a long-lived GPU worker process, isolated from the engine. \\
Persistence & SQLite (WAL, numbered migrations) and a content-addressed asset store. \\
Incremental work & Every pipeline step has a cache key; re-sync resubmits everything and finished steps hit the cache. \\
\bottomrule
\end{tabularx}

\chapter{Building block view}
\begin{tabularx}{\linewidth}{@{}lY@{}}
\toprule
\textbf{Module} & \textbf{Responsibility} \\
\midrule
\texttt{domain}, \texttt{config} & Data models, paths under \texttt{\%LOCALAPPDATA\%}, secrets in the Windows credential store. \\
\texttt{store} & SQLite access, migrations, asset store. \\
\texttt{sources} & Local files, watched folders, playlist-file import, Creative Commons catalogues. \\
\texttt{library} & Collections, matching of imported entries to local files, re-sync, exports. \\
\texttt{audio} & Decoding and encoding via ffmpeg, mixer with \gls{residual}, waveform peaks. \\
\texttt{separation} & Model catalogue and installer, recipe graphs, GPU worker supervisor, hardware detection. \\
\texttt{pipeline} & Steps, job runner, cancellation, event bus. \\
\texttt{api}, \texttt{cli} & FastAPI adapter and Typer command line. \\
\bottomrule
\end{tabularx}

\chapter{Runtime view}
\todoin{M3}

\chapter{Deployment view}
\todoin{M9}

\chapter{Cross-cutting concepts}
\section{Local API security}
The engine binds to 127.0.0.1 on a random port. The shell passes a 32-byte token
through an environment variable; every request must carry it, use the exact
\texttt{Host} header and an allowed \texttt{Origin}. WebSocket clients
authenticate with their first message.

\section{Hardware adaptation}
The engine detects GPU vendor, compute capability, free \gls{vram}, RAM, cores,
power source and free disk. On \gls{oom} it retries with half precision and
smaller chunks before asking the user to fall back to the CPU. It also watches
throughput to detect the driver's system-memory fallback, which slows inference
instead of failing.

\section{DAW coexistence}
Separation runs at reduced priority with an efficiency hint and a thread cap. If a
\gls{daw} is running, a low-impact mode with smaller chunks and pauses is offered.

\chapter{Architecture decisions}
Decisions are recorded as \glspl{adr}. Planned: ADR-001 separation engine and
presets, ADR-002 \gls{vram} tiers and fallback (M1); ADR-004 transport and
security, ADR-005 process model, ADR-006 bootstrap, ADR-007 audio preview (M2).

\chapter{Compliance}\label{ch:compliance}
\cmnote{This chapter summarises a risk analysis carried out on 2026-11-02. It is
not legal advice. Jurisdiction-dependent points are marked as such.}

\section{Acquisition of audio}
\begin{itemize}
  \item Spanish private-copy exception (TRLPI art.~31.2): requires a lawful source
        and respect for the conditions of access; it cannot be invoked against
        technological protection measures (art.~197.4). The CJEU excluded
        unlawful sources from private copying (ACI Adam, C-435/12).
  \item German courts held that a major video platform's stream protection is an
        effective \gls{tpm} (Hamburg Higher Regional Court, 5~U~54/23, 2024).
        Distributing a tool that circumvents it may create liability under
        TRLPI art.~196.2 or 17~U.S.C.~\S1201.
  \item A major streaming service's developer policy and terms forbid using its
        metadata to obtain the same music from another service; in October 2026
        a project following that pattern was removed from GitHub.
\end{itemize}
\textbf{Decision:} \CM{} does not acquire audio from streaming platforms and does
not use their \glspl{api} or web pages (constraints CON-01 and CON-02).

\section{Stems and publication}
Separating lawfully obtained audio for private practice carries very low risk.
Publishing a \gls{backingtrack} made from an original recording requires rights to
both the recording and the composition; the application warns users on export.

\section{Model weights}
Most community separation models have no explicit licence, a non-commercial
licence or a copyleft licence, and the provenance of some is unclear. \CM{}
therefore ships no default model, never hosts weights, shows licence class and
provenance, and records the exact licence text each user accepted (CON-03).

\section{Trademarks}
Third-party names appear only as plain text describing formats or compatibility
(CON-05).

\chapter{AI model evaluation}
\section{Survey (2026-11)}
Signal-to-distortion ratios (\gls{sdr}, higher is better) from the public MVSep
multisong leaderboard:

\begin{tabularx}{\linewidth}{@{}lYll@{}}
\toprule
\textbf{Model} & \textbf{SDR per stem (dB)} & \textbf{Licence} & \textbf{Loader} \\
\midrule
BS-Roformer-SW (6 stems) & vocals 11.30, bass 14.62, drums 14.11, guitar 9.01, piano 7.80, other 8.72 & none / unclear & audio-separator \\
unwa Resurrection (vocals) & 11.34 & none & audio-separator \\
unwa Leap Xe (vocals) & 11.79 & none & MSST \\
Kim Mel-Band (vocals) & about 11.0 & MIT (unconfirmed) & audio-separator \\
HT Demucs / 6s & drums 10.9, bass 11.8, guitar 5.25 & disputed & audio-separator \\
Previous v0.1 model (becruily vocals) & about 10.55 & none & audio-separator \\
\bottomrule
\end{tabularx}

\section{Estimated performance per hardware tier}
Estimates for a 4-minute song (\(\pm 50\%\)), to be replaced by M1 measurements:

\begin{tabularx}{\linewidth}{@{}lYYY@{}}
\toprule
\textbf{Tier} & \textbf{BS-Roformer-SW} & \textbf{Vocal RoFormer} & \textbf{HT Demucs} \\
\midrule
CPU only & 20--35 min & 20--40 min & 2--4 min \\
4 GB VRAM & 1.5--4 min & 1--3 min & 20--40 s \\
6 GB VRAM & 1--2.5 min & 45 s--2 min & 15--30 s \\
8 GB VRAM & 40--90 s & 35--80 s & 10--20 s \\
12 GB+ VRAM & 20--60 s & 20--40 s & 6--12 s \\
\bottomrule
\end{tabularx}

\section{Benchmarks}
\todoin{M1}

\printnoidxglossaries
\end{document}
